% TOP_PROBLEM: {"id": 5500060, "title": "Transforming Polygons via Vertex-Centroid Moves", "category_id": 6, "category": {"id": 6, "name": "geometry", "display_name": "Geometry", "description": "Euclidean and non-Euclidean geometry, geometric structures.", "slug": "geometry", "order_index": 6, "created_at": "2026-07-31T15:26:25.671Z"}, "status": "open"}
% FIRST_POSTED: null
% ENTRY_KIND: statement_only
% Imported from: batch07.tex; UnsolvedMath ID 5500060
\documentclass[11pt]{article}
\usepackage{amsmath,amssymb,mathtools}
\usepackage{fontspec,unicode-math}
\setmainfont{Latin Modern Roman}
\setsansfont{Latin Modern Sans}
\setmathfont{Latin Modern Math}
\usepackage{xurl,hyperref}
\newcommand{\sref}[2]{\hyperlink{source:#1:#2}{[S#2]}}
\newcommand{\cataloguescope}[1]{\paragraph{Mathematical question.}#1}
\begin{document}
% UnsolvedMath ID 5500060; source catalogue code AMR-054-0060
\setcounter{section}{5500059}
\section{Transforming Polygons via Vertex-Centroid Moves}\label{5500060}
{\small\sffamily UnsolvedMath ID 5500060 · Geometry}
\subsection{Definitions and mathematical statement}

\paragraph{Shared notation.}$\mathbb N=\{1,2,\ldots\}$, and $\mathbb Z,\mathbb Q,\mathbb R,\mathbb C$ denote the integers, rationals, reals and complex numbers. Any other conventions are specified locally.


For vertices $v_1,\ldots,v_n\in\mathbb R^2$, the vertex centroid is $m=n^{-1}\sum_i v_i$. In one vertex-centroid move, a single vertex $v_i$ is translated along the line through $v_i$ and the current $m$, with all other vertices fixed. Recompute the centroid after each move. Vertices can be moved in any order and can be moved repeatedly. The same vertex-centroid definition applies in $\mathbb R^d$. A scaling target uses a prescribed factor $s>0$; a rotation target applies a prescribed Euclidean rotation. When a polygon bounds a region $P$ of positive area, its area centroid is $m_{\rm area}=\operatorname{area}(P)^{-1}\int_P x\,dx$. In the area-centroid variant, the moving vertex moves along the instantaneous line through it and the current area centroid, which can itself leave the initial line during the move.
\paragraph{Statement recorded in the supplied source.}
Can every polygon be transformed exactly into a regular polygon by a finite sequence of vertex-centroid moves? Investigate also the associated transformation goals of scaling by a prescribed $s>0$ or rotating the polygon, and vertex-centroid transformations in arbitrary dimensions. Does the regularization question hold with the area centroid in place of the vertex centroid? \sref{5500060}{2}
\subsection{Short English statement}
Can every polygon be transformed exactly into a regular polygon by a finite sequence of vertex-centroid moves? Investigate also the associated transformation goals of scaling by a prescribed $s>0$ or rotating the polygon, and vertex-centroid transformations in arbitrary dimensions. Does the regularization question hold with the area centroid in place of the vertex centroid?
\subsection{Sources}
\begin{description}
\item[\hypertarget{source:5500060:1}{S1}] ULAM.AI and the UnsolvedMath Contributors, UnsolvedMath: A Curated Collection of Open Mathematics Problems, record 5500060 (AMR-054-0060). CC BY 4.0. \url{https://huggingface.co/datasets/ulamai/UnsolvedMath}.
\item[\hypertarget{source:5500060:2}{S2}] Erik D. Demaine, Joseph S. B. Mitchell and Joseph O\textquotesingle Rourke (editors), \emph{The Open Problems Project}, maintained original source, Problem 60, Statement and complete Partial/Related Results. \url{https://github.com/edemaine/topp/blob/main/Problems/P.000060}.
\end{description}
\label{end:5500060}
\end{document}
